% 07-serving
\section{Serving and matching}\label{sec:serving}

The offline path leaves one compact profile per URL. At bid time the bidder
has to find that profile in one round trip, never wait for a crawl, decide
whether a deal or a creative may run there, say why, and rank what is left. We
split the work by language. The OCaml pipeline writes the profiles. The
bidder, written in Go, reads them through a small SDK. The rules therefore
exist twice, and both sides are checked against one set of fixtures that the
OCaml side writes.

\subsection{One key per page}\label{sec:serving-key}

A profile is found by its URL, so both sides must spell a URL the same way. Let
$\nu(u)$ be the bidder's own URL normalisation, the one it already applies
before its other per-URL lookups. The key of page $u$ is
\begin{equation}\label{eq:serving-key}
\kappa(u) = \texttt{contextual:} \,\|\, \operatorname{hex}\bigl(\operatorname{SHA\text{-}256}(\nu(u))\bigr).
\end{equation}
The OCaml side ports $\nu$ byte for byte from Go's \code{net/url} parser. It
drops the fragment, the scheme and any user information. It unescapes the
path, removes a leading \code{www.}\ from the host and one trailing slash from
the path, takes a fixed list of tracking parameters such as \code{utm\_source}
and \code{fbclid} out of the query, and removes a trailing \code{/amp}. The
host's case and port are kept. So
\code{https://www.example.com/news/article/amp?utm\_source=x\&fbclid=y\#top}
becomes \code{example.com/news/article}. A URL that Go refuses is refused
here too, so no profile is written under a key the bidder cannot compute. Any
difference would be a silent miss, so the port is tested on 111 cases written
by a Go implementation, 16 of them refused, and the Go SDK must recompute the
other 95 keys exactly.

\subsection{The serving profile}\label{sec:serving-profile}

The value under $\kappa(u)$ is a protobuf message, \code{PageProfile}, with a
30-day expiry (schema in \cref{app:wire}). It holds the \preset{rich} level's
labels, at most 20, highest score first; at most 20 keywords; the fetch time;
the preset's name; and the pipeline version, a hash of the taxonomy, lexicon,
models, calibration, thresholds and \tier{4} model, for logs and audits. The labels,
keywords, fetch time and preset take 252 bytes on average. For slot context
(\cref{sec:boundary}) the message also carries up to 64 labelled blocks, each
as its SimHash, its \code{main} flag, up to 8 labels and its estimated
position, and the page's ad containers. No block text is stored. The lab's full
profile, with every level, all evidence and the segments' text, is a separate
value of 9.3\,KB per page after zstd compression, which the bidder never
reads.

The labels obey one invariant, which saves the bidder a copy of the taxonomy.

\begin{definition}[Ancestor-closed profile]\label{def:serving-closed}
A label list $L$ with scores $s$ is \emph{ancestor-closed} when for every
$v \in L$ and every $u \in \operatorname{anc}(v)$, also $u \in L$ and
$s(u) \ge s(v)$.
\end{definition}

The cascade adds ancestors to every label (\cref{sec:cascade}), so served
profiles have this property. A rule on a node then needs no tree: ``Sports or
anything under it scores at least 0.5'' holds exactly when the id of Sports
itself is in the list with a score of at least 0.5. Checking one category
in the Go SDK is a scan of at most 20 labels.

\subsection{Backoff: page, section, site}\label{sec:serving-backoff}

A new article has no profile until the crawler reaches it, and most requests
miss while a crawl is young. So a page that misses inherits a profile from
where it lives. Write a normalised URL as $h/p_1/p_2/\dots$, with host $h$ and
the query left out. Its top section is $h/p_1$ (none when $p_1$ is empty) and
its site is $h$. Both have their own key,
\begin{equation}\label{eq:serving-section-key}
\kappa_s(q) = \texttt{contextual:section:} \,\|\, \operatorname{hex}\bigl(\operatorname{SHA\text{-}256}(q)\bigr),
\end{equation}
in a separate namespace, because a section's top page normalises to the same
string as the section. The bidder tries $\kappa(u)$, then the section's key,
then the site's (\cref{fig:serving-sequence}). We stop at the top section on
purpose: the bidder's store client may have no multi-key read, so each level
costs one more round trip on a miss.

A section's profile is the mean of its pages. Let $D_q$ be the pages saved
under prefix $q$, $n_q = |D_q|$, and $s_d(v)$ the served score of $v$ on page
$d$, with $s_d(v) = 0$ when $v$ is not among its served labels. The store keeps
the running sums $S_q(v) = \sum_{d \in D_q} s_d(v)$ and the count $n_q$. When
the crawler saves page $d$ with new scores $s'_d$, and $s_d$ are the scores it
added on its last crawl (all $0$ for a new page),
\begin{align}
S_q(v) &\leftarrow S_q(v) + s'_d(v) - s_d(v) \quad\text{for every } v \text{ in either list},\label{eq:serving-sums}\\
n_q &\leftarrow n_q + [\,d \text{ is new}\,],\nonumber\\
\bar{s}_q(v) &= S_q(v) / n_q, \quad\text{served once } n_q \ge 5.\label{eq:serving-mean}
\end{align}
A recrawled page therefore takes back what it added before, and the mean stays
the mean of each page's latest scores. One Lua script applies
\cref{eq:serving-sums} to the section and the site in a single atomic step,
deletes sums that fall below $10^{-6}$ and returns the hashes. The crawler then
writes each prefix with at least 5 pages as a \code{PageProfile}: its 20
highest means, no keywords, under $\kappa_s(q)$ with the same 30-day expiry.
The previous scores come from the page's summary in the store, so no second
copy of the profile is read.

\begin{proposition}\label{prop:serving-closed}
If every page's served labels are ancestor-closed, so are the section's means
in \cref{eq:serving-mean}, before the cut to the 20 highest.
\end{proposition}
\begin{proof}
Take $u \in \operatorname{anc}(v)$. On every page $s_d(u) \ge s_d(v)$: when $v$
is served this is \cref{def:serving-closed}, and otherwise $s_d(v) = 0$. Sum
over $D_q$ and divide by $n_q$.
\end{proof}

So \code{Rule.Check} reads a section profile as it reads a page's; the cut to
20 can only break this on a tie at the boundary. On the golden crawl, 383 pages
give 37 sections and sites with at least 5 pages. On the Pi cluster, adding
3,236 stored pages took 37 seconds and produced 1,318 prefixes, 242 of which
reached 5 pages after a run with synthetic traffic. Hubs such as home pages,
section fronts and live blogs change on almost every fetch, so the crawler
deletes their page key on purpose and the bidder backs off to their section,
which they keep feeding. Two limits remain. A page counts in its section until
it is recrawled, however old; a window or a decay would fix that if sections
drift. And URLs that start with a date give a year as their ``section''.

\begin{figure}[t]
\centering
\begin{tikzpicture}[
  head/.style={box,minimum width=2.6cm,minimum height=0.75cm,font=\footnotesize},
  msg/.style={font=\footnotesize,above,inner sep=1.5pt,fill=white},
  act/.style={draw=ink,fill=soft,minimum width=0.22cm,inner sep=0pt},
  life/.style={draw=ink!40,dashed}]
  \foreach \x/\n/\t in {0/bid/{bid path\\(Go SDK)},4.1/st/{shared store},8.2/wk/{crawl-request\\workers},12.3/cr/{crawler}} {
    \node[head] (\n) at (\x,0) {\t};
    \draw[life] (\x,-0.4) -- (\x,-8.6);
  }
  % Miss on the page key.
  \draw[arrow] (0,-1.0) -- node[msg] {\code{GET} $\kappa(u)$} (4.1,-1.0);
  \draw[dashedarrow] (4.1,-1.45) -- node[msg] {nothing} (0,-1.45);
  % Section hit.
  \draw[arrow] (0,-2.0) -- node[msg] {\code{GET} $\kappa_s(h/p_1)$} (4.1,-2.0);
  \draw[dashedarrow] (4.1,-2.45) -- node[msg] {mean profile, $n_q \ge 5$} (0,-2.45);
  \node[note,anchor=west,align=left] at (4.25,-2.22) {on a miss, one more\\\code{GET} $\kappa_s(h)$};
  % Local work on the bid path.
  \node[act,minimum height=0.55cm] at (0,-3.05) {};
  \node[font=\footnotesize,anchor=west,align=left] at (0.25,-3.05) {decode, \code{Rule.Check};\\level \code{Section}};
  % Crawl request, never blocking.
  \draw[arrow] (0,-3.8) -- node[msg,pos=0.72] {\code{Request}$(\nu(u), u)$} (8.2,-3.8);
  \draw[dashedarrow] (8.2,-4.25) -- node[msg,pos=0.72] {queued, at once} (0,-4.25);
  \node[note,anchor=west] at (0.15,-4.75) {the bid goes on};
  \node[act,minimum height=0.5cm] at (8.2,-5.0) {};
  \draw[arrow] (8.2,-5.0) -- node[msg] {publish $u$} (12.3,-5.0);
  % The crawler, later.
  \node[act,minimum height=1.3cm] at (12.3,-6.0) {};
  \node[font=\footnotesize,anchor=east,align=right] at (12.1,-5.85) {fetch politely, run\\the cascade (later)};
  \draw[arrow] (12.3,-6.6) -- node[msg] {\code{SET} $\kappa(u)$, 30 days; section sums} (4.1,-6.6);
  % Next request for the same page.
  \draw[arrow] (0,-7.55) -- node[msg] {\code{GET} $\kappa(u)$} (4.1,-7.55);
  \draw[dashedarrow] (4.1,-8.0) -- node[msg] {page profile} (0,-8.0);
  \node[font=\footnotesize,anchor=west] at (0.15,-8.35) {level \code{Page}};
  % Braces.
  \draw[decorate,decoration={brace,amplitude=4pt,mirror},ink!60] (-0.25,-0.85) -- (-0.25,-4.4)
    node[midway,left=5pt,font=\footnotesize,align=right,text=ink] {one request:\\up to 3 \code{GET}s,\\no waiting};
  \draw[decorate,decoration={brace,amplitude=4pt,mirror},ink!60] (-0.25,-7.4) -- (-0.25,-8.5)
    node[midway,left=5pt,font=\footnotesize,align=right,text=ink] {a later\\request};
\end{tikzpicture}
\caption{A bid-time lookup for a page nobody has crawled yet. The page key
  misses, the top section's mean profile answers, and the deal's rule runs on
  it locally. The crawl request returns at once; a worker publishes it off the
  bid path, and a later request for the same page finds its own profile.}
\label{fig:serving-sequence}
\end{figure}

\subsection{Crawl requests on a miss}\label{sec:serving-crawl}

Every request that does not find its page's own profile asks the crawler for
it. The call must never slow the bid, so \code{CrawlRequests.Request} does only
a map lookup and a channel send. It first marks the normalised URL as asked;
a URL asked before returns \emph{already asked}. It then offers the page to a
bounded queue and returns \emph{queued}, or \emph{dropped} when the queue is
full or closed. A fixed pool of workers takes pages off the queue and
publishes each URL, one message per page, to a topic the crawler subscribes
to. A page that is dropped, or whose publish fails, is forgotten, so the next
miss asks again. The set of asked pages starts over at 100,000 entries; a
repeat then costs one message, which the crawler drops. The path from the SDK to the crawler was checked end to
end on an emulator of the message service, with its real client library. The
topic, the subscription and the permissions belong
to the bidder's infrastructure and are not set up yet.

\subsection{Hard rules, then a soft score}\label{sec:serving-rules}

Matching runs in two stages (\cref{tab:serving-pipeline}). Hard rules come
first, in a fixed order, and each refusal comes with a sentence a trader can
read. A label counts once its score is at least 0.5, a node counts together
with its descendants, and the reason names the deepest node that matched.
Blocks are none-of: blocking Sports also blocks Soccer. Suitability refuses a
creative that avoids Sensitive Topics when any node under that branch is on
the page. Targets are any-of, and a descendant counts, so targeting Food \&
Drink matches Cooking. When a campaign under-delivers, the
logged refusals say which rule starves it.

\begin{table}[t]
\centering
\caption{The matcher's stages (\code{Matcher.refusal} and \code{Matcher.score}
  in \code{src/lib/matcher.ml}), in order, with the reasons the test suite
  prints. The Go SDK ports the three hard rules as \code{Rule.Check}; the soft
  score stays in OCaml.}
\label{tab:serving-pipeline}
\small
\begin{tabularx}{\linewidth}{@{}l >{\raggedright\arraybackslash}p{4.1cm} >{\raggedright\arraybackslash}X l@{}}
\toprule
Stage & Passes when & Reason printed when it decides & Go verdict\\
\midrule
1. blocks & no blocked node on the page & advertiser blocks Sports; page is Sports > Soccer (0.92) & blocked\\
2. suitability & the creative accepts Sensitive Topics, or none is on the page & brand suitability: page is Sensitive Topics > Crime \& Harmful Acts\dots\ (0.89) & unsuitable\\
3. targets & no targets, or one is on the page & no targeted category on page & off-target\\
4. soft score & stages 1--3 passed & topic Sports > Soccer $\sim$ Sports > Soccer (0.83); keywords: premier league & not ported\\
\bottomrule
\end{tabularx}
\end{table}

Creatives that pass get a score from three named features. The
first is topical. Let page $d$ carry labels $u$ with scores $s_u$ and creative
$c$ carry labels $w$ with scores $t_w$, both in the same \iab{} label space. With
$\operatorname{depth}(v) = |\operatorname{anc}(v)|$, so that a root has depth
1, the Wu--Palmer similarity of two nodes is
\begin{equation}\label{eq:serving-wup}
\operatorname{wp}(u, w) = \frac{2\,\operatorname{depth}(\operatorname{lca}(u, w))}{\operatorname{depth}(u) + \operatorname{depth}(w)}
 = \frac{2\,|\operatorname{anc}(u) \cap \operatorname{anc}(w)|}{|\operatorname{anc}(u)| + |\operatorname{anc}(w)|},
\end{equation}
and $0$ when $u$ and $w$ sit under different roots (\cref{fig:serving-wup}).
It is 1 for the same node, $2/3$ for a root and its child, and $1/2$ for two
children of a root. The topic feature takes the best pair,
\begin{equation}\label{eq:serving-topic}
\operatorname{topic}(d, c) = \max_{u, w}\; s_u\, t_w\, \operatorname{wp}(u, w),
\end{equation}
and reports that pair as its reason. This is the taxonomy score of
\citet{broder2007semantic}, which backs off through the lowest common ancestor
in the same way. The keyword feature is the share of the creative's keywords
$K_c$ found among the page's $K_d$, after both pass through the same text
normaliser: $\operatorname{kw}(d, c) = |K_c \cap K_d| / |K_c|$, and $0$ when
$K_c$ is empty. Requiring the ad's keywords on the page lifted average
precision by 60\% over plain cosine matching in \citet{ribeironeto2005impedance}. The lift feature
is $\operatorname{lift}(d, c) = \max_u s_u\, A(\pi_c, u)$, where $\pi_c$ is the
creative's product category and $A$ a table of learned affinities between
product and page categories. Because $A$ is keyed by product, a new creative
inherits what its category has already taught, and beer can score on a sports
page without any topical overlap. The score is
\begin{equation}\label{eq:serving-score}
\operatorname{score}(d, c) = 0.6\,\operatorname{topic}(d, c) + 0.25\,\operatorname{kw}(d, c) + 0.15\,\operatorname{lift}(d, c).
\end{equation}
On a football match report the test suite gives boots 0.621, from the pair
Soccer $\sim$ Soccer (0.83) and the keyword ``premier league''; beer 0.083,
from a test lift of Beer on Sports; an electric car blocked by its
advertiser's block on Sports; and cookware off-target.

\begin{figure}[t]
\centering
\begin{tikzpicture}[
  tn/.style={box,font=\footnotesize,minimum width=1.9cm,minimum height=0.62cm,inner sep=3pt},
  hl/.style={tn,draw=accent,line width=1.1pt,fill=accent!12},
  dk/.style={tn,draw=ink,line width=1.1pt},
  tag/.style={font=\scriptsize\bfseries,inner sep=1pt},
  edge/.style={draw=ink!60}]
  \foreach \y/\t in {0/depth 1,-1.35/depth 2,-2.7/depth 3} {
    \draw[ink!15,dashed] (-2.0,\y) -- (12.2,\y);
    \node[font=\scriptsize,text=ink!60,anchor=east] at (-2.05,\y) {\t};
  }
  \node[tn] (sp) at (3.2,0) {Sports};
  \node[tn,draw=accent,line width=1.1pt] (cs) at (0.9,-1.35) {College Sports};
  \node[dk] (so) at (3.75,-1.35) {Soccer};
  \node[tn] (ru) at (6.1,-1.35) {Rugby};
  \node[hl] (cf) at (-0.55,-2.7) {College Football};
  \node[hl] (cb) at (2.45,-2.7) {College Basketball};
  \node[tn] (rl) at (6.1,-2.7) {Rugby Union};
  \node[tn] (fd) at (9.85,0) {Food \& Drink};
  \node[tn] (co) at (8.75,-1.35) {Cooking};
  \node[tn] (ab) at (11.05,-1.35) {Alc.\ Beverages};
  \foreach \a/\b in {sp/cs,sp/so,sp/ru,cs/cf,cs/cb,ru/rl,fd/co,fd/ab} \draw[edge] (\a.south) -- (\b.north);
  \node[tag,text=accent,anchor=south east] at (cs.north west) {LCA};
  \node[tag,text=ink,anchor=south east] at (sp.north west) {LCA};
  \node[font=\footnotesize,align=left,anchor=north west] at (-2.0,-3.3) {%
    \textcolor{accent}{$\operatorname{wp}(\text{College Football}, \text{College Basketball}) = \tfrac{2\cdot 2}{3+3} = 0.67$},
      lowest common ancestor College Sports\\[3pt]
    $\operatorname{wp}(\text{College Football}, \text{Soccer}) = \tfrac{2\cdot 1}{3+2} = 0.40$,
      lowest common ancestor Sports\\[3pt]
    $\operatorname{wp}(\text{Soccer}, \text{Cooking}) = 0$, no common ancestor};
\end{tikzpicture}
\caption{Wu--Palmer similarity on two branches of \iab. Two leaves under the
  same parent (filled blue) meet at College Sports and score 0.67. College
  Football and Soccer (dark outline) meet only at the root and score 0.40.
  Nodes under different roots score 0, so a topic match never crosses from
  Sports to Food \& Drink.}
\label{fig:serving-wup}
\end{figure}

The weights 0.6, 0.25 and 0.15 are set by hand. They are a placeholder for a
learned model, marked as a shortcut in the code, and we make no claim that
they are right. The plan is a logistic regression on engagement over the same
named features with format and device crosses, as
\citet{chakrabarti2008contextual} did for clicks, so that the weights stay
readable as logits. A contextual bandit over the same features
\citep{li2010contextual} would then explore among eligible creatives, with
propensities logged for offline evaluation, and a floor below which the bidder
does not force relevance and falls back to its non-contextual bid
\citep{broder2008swing}. The lift table is empty today, so lift is 0 in
every score today. The creative side is also thin: a creative comes from a text
prompt labelled by \tier{1} and \tier{2}, not yet from its ad copy and landing
page, and product categories are not loaded.

\subsection{Curated deals: placement list or context}\label{sec:serving-deals}

The first use is curated deals. A deal takes a request whose placement is on
its list, \emph{or} whose page matches its context. If deal filters are ANDed,
a separate context filter could only shrink supply, so the OR belongs inside
the placement check:
\begin{lstlisting}
func placementOK(deal *Deal, req *BidRequest) bool {
    if deal.Placements.Contains(req.PlacementID) {
        return true
    }
    return deal.ContextRule != nil && req.ContextualLevel != contextual.NotFound &&
        deal.ContextRule.Check(&req.Contextual).Verdict == contextual.Eligible
}
\end{lstlisting}
A deal's context is a \code{Rule}: its targets, blocks, Sensitive Topics choice
and a minimum score, 0.5 by default. A page with no profile at any level is not
eligible, as keyword targeting treats a request without keywords. When the
request names the blocks around its slot, the rule can run on the slot's
context from \code{Profile.AtSlot} instead of the page's
(\cref{sec:boundary}). The bidder integration is written as a guide, not
deployed. Its first step is a shadow mode that only counts, per deal, the
requests context would add. A per-deal minimum score that falls when a deal is
behind pace and rises when it is ahead is planned, as are per-bid logs of each
decision.

\subsection{An allocation-free reader}\label{sec:serving-reader}

A bidder decodes a profile on every request, so the decoder must not allocate.
\code{Profile.Decode} walks the protobuf wire format by hand. It reuses the
profile's slices, so a pooled profile decodes in place. Label ids are not
copied: they point into the input buffer, which must not change while the
profile is in use. Every block's labels share one backing array. Fields the
decoder does not know are skipped, so a writer can add fields without breaking
deployed readers. The tests assert zero allocations for decoding a profile
with blocks and for a decode followed by \code{Rule.Check}.
\Cref{tab:serving-sdk} lists the calls, their cost and the fixtures that pin
each one to the OCaml side.

\begin{table}[t]
\centering
\caption{The Go SDK's calls (\code{sdks/go/contextual}). Costs are Go
  benchmarks on an Apple M4. Fixtures are written by the OCaml side, so both
  implementations answer the same cases.}
\label{tab:serving-sdk}
\small
\begin{tabularx}{\linewidth}{@{}l >{\raggedright\arraybackslash}X l l@{}}
\toprule
Call & What it does & Cost per call & Checked against\\
\midrule
\code{Key} & the page's key, \cref{eq:serving-key} & 83\,ns, 2 allocations & 95 URLs\\
\code{Sections}, \code{SectionKey} & prefixes and their keys, \cref{eq:serving-section-key} & & the same 95 URLs\\
\code{LookupWithBackoff} & page, then section, then site & one \code{GET} per level & unit tests\\
\code{Profile.Decode} & protobuf into a reused profile & 83\,ns, no allocation & every fixture\\
\code{Rule.Check} & a deal's rule and its reason & 9\,ns, no allocation & 3,000 profiles\\
\code{Profile.AtSlot} & the slot's context (\cref{sec:boundary}) & 1.8\,\textmu s, 23 allocations & 2,000 cases\\
\code{Profile.SlotOfGPID} & the slot's blocks from a GPID & & 1,000 cases\\
\code{CrawlRequests.Request} & a crawl request that never blocks & a map lookup, a send & unit tests\\
\bottomrule
\end{tabularx}
\end{table}

The 3,000 rule fixtures are profiles built as the cascade builds them, with
ancestors added and scores rounded to 32-bit floats as the wire carries them,
each paired with a random deal rule and the OCaml matcher's verdict. The 2,000
slot fixtures pair a profile with the \code{ctx} of a bid request. Most of its
hashes copy one of the profile's blocks with up to 5 bits flipped and the rest
are random, so both sides of the 3-bit threshold are exercised. \code{AtSlot} must match \code{Slot.context}
to within $10^{-5}$ on every label. The fixtures are regenerated after any
change to the matcher, the label order or the profile format.
